% UBT-AI-PROVENANCE-BEGIN
% schema: ubt-ai-provenance/v1
% tier: C_working
% ai_assistance: disclosed
% human_review: risk-based
% editorial_responsibility: Ing. David Jaroš
% policy: ../../../AI_PROVENANCE.md
% notice: Working material; exhaustive human review is not claimed.
% UBT-AI-PROVENANCE-END

\chapter*{Předmluva: jak tuto učebnici číst}
\addcontentsline{toc}{chapter}{Předmluva}

Tato kniha je pracovní učebnicový a referenční text programu Unified
Biquaternion Theory (UBT). Jejím prvním úkolem je podat matematiku natolik
srozumitelně a souvisle, aby bylo možné později znovu rekonstruovat vazby mezi
jednotlivými větvemi UBT bez spoléhání na paměť nebo staré chaty. Proto jsou v
textu záměrně odvozovány i klasické matematické výsledky, pokud jsou potřebné
pro velký obraz UBT, přestože samotná matematika může být desítky či stovky let
známá.

Učebnice rozlišuje tři vrstvy:
\begin{enumerate}
  \item\textbf{klasické základy}: standardní matematika a fyzika s dostatečným
        odvozením, aby byl text co nejvíce soběstačný;
  \item\textbf{propojení s UBT}: kontrolované převody těchto standardních struktur
        do UBT notace, omezených modelů nebo konkrétních hypotéz;
  \item\textbf{tvrzení UBT a otevřené mezery}: výroky závislé na kanonickém poli
        UBT, jeho akci, dynamice nebo empirickém programu.
\end{enumerate}

Klasická identita se nestává novým výsledkem jen proto, že ji UBT používá.
Stejně tak ale skutečnost, že jde o známou matematiku, není důvod ji z učebnice
vyhodit: výzkumná hodnota může spočívat v tom, jak jsou známé struktury složeny
dohromady a jakou novou otázku pro UBT jejich spojení přesně formuluje.
Kanonické dokumenty tvrzení proto zůstávají konzervativní, zatímco učebnice může
být široká a didaktická, pokud jsou stavy tvrzení vždy zřetelné.

Živá učebnice nesmí přebírat aktivní teorii z \texttt{ARCHIVE/}. Aktuální
stav se nejprve čte ze směrodatných registrů, zejména \texttt{STATUS.md},
\texttt{WHAT\_IS\_PROVED.md} a \texttt{CLAIMS.yaml}; teprve potom se používá
aktivní matematika z \texttt{canonical/}. Standardní matematiku lze v učebnici
odvodit přímo. Historické soubory slouží pouze při popisu historie myšlenky.
Migrační mapa je v \texttt{docs/textbook/SOURCE\_MAP.md}.

Primární studentská edice je česká. Anglická verze může být později vedena jako
paralelní edice se stejnou matematickou strukturou a stejnými statusy tvrzení;
nové kapitoly se nemají psát jako náhodná směs češtiny a angličtiny.

Každý aktuální zdrojový soubor zároveň používá strojově čitelná pravidla
repozitáře pro provenienci AI. Tato učebnice je pracovním materiálem úrovně C,
dokud lidský editor výslovně nezmění úroveň ve směrodatné mapě provenience;
agent AI nesmí
provádět vlastní povýšení nebo atestaci.

Odvození v článcích a části učebnice důležité pro teorémy se mají nezávisle
kontrolovat podle \texttt{docs/DERIVATION\_VERIFICATION\_POLICY.md}. Pro
formalizovatelnou přesnou matematiku je preferovaný Lean; vedle něj se používají
nezávislé systémy počítačové algebry (CAS) a numerické kontroly. Úspěšný výpočet v nástroji však ověřuje jen
to, co bylo do nástroje skutečně zakódováno, a sám o sobě nemění fyzikální
stav tvrzení UBT.
